
Para ver os dois métodos funcionando com números, considere a recorrência
\begin{align*}
  f(n) = 2 f(n-1) + 3 f(n-2), \\
  f(0) = 1, \\
  f(1) = 4,
\end{align*}
ou seja, $\alpha = 2$, $\beta = 3$, $c = 1$ e $d = 4$. Seus primeiros termos são $1, 4, 11, 34, 101, 304, 911, \dots$

Os valores foram escolhidos para que as contas fiquem limpas: $\beta \neq 0$ (a recorrência é realmente de ordem 2), $\alpha^2 + 4\beta = 16 \neq 0$ (dois autovalores distintos) e, como $16$ é um quadrado perfeito, os autovalores são inteiros. Além disso, $d \neq c \lambda_1$ e $d \neq c \lambda_2$, de modo que nenhuma coordenada de $\bar{f}^{(0)}$ se anula e as duas componentes aparecem na resposta.

\section{Pela diagonalização}

A matriz do sistema é
\[
  A =
  \begin{bmatrix}
    0 & 1 \\
    3 & 2
  \end{bmatrix},
  \quad
  f^{(0)} =
  \begin{bmatrix}
    1 \\
    4
  \end{bmatrix}.
\]

Calculando os autovalores:
\begin{calcproof}[.]
  \start{\det{A - \lambda I} = 0}
  \step[\calciff]{Definição ($A$) e Álgebra}{
    \det{\begin{bmatrix}
        - \lambda & 1 \\
        3 & 2 - \lambda
      \end{bmatrix}} = 0
  }
  \step[\calciff]{Cálculo do determinante $2 \times 2$ e Álgebra}{
    \lambda^2 - 2 \lambda - 3 = 0
  }
  \step[\calciff]{Fatoração}{
    (\lambda - 3)(\lambda + 1) = 0
  }
  \step[\calciff]{Produto nulo}{
    \lambda_1 = 3 \ \lor \ \lambda_2 = -1
  }
\end{calcproof}

Pelo capítulo da diagonalização, os autovetores são $v_i = \begin{bmatrix} 1 \\ \lambda_i \end{bmatrix}$, então
\[
  V =
  \begin{bmatrix}
    1 & 1 \\
    3 & -1
  \end{bmatrix},
  \quad
  V^{-1} = \dfrac{1}{\lambda_2 - \lambda_1}
  \begin{bmatrix}
    \lambda_2 & -1 \\
    -\lambda_1 & 1
  \end{bmatrix}
  = \dfrac{1}{-4}
  \begin{bmatrix}
    -1 & -1 \\
    -3 & 1
  \end{bmatrix}
  = \dfrac{1}{4}
  \begin{bmatrix}
    1 & 1 \\
    3 & -1
  \end{bmatrix}.
\]

As coordenadas do estado inicial na base de autovetores são
\[
  \bar{f}^{(0)} = V^{-1} f^{(0)}
  = \dfrac{1}{4}
  \begin{bmatrix}
    1 & 1 \\
    3 & -1
  \end{bmatrix}
  \begin{bmatrix}
    1 \\
    4
  \end{bmatrix}
  = \dfrac{1}{4}
  \begin{bmatrix}
    5 \\
    -1
  \end{bmatrix},
\]
isto é,
\[
  \bar{f}^{(0)}_1 = \dfrac{5}{4} \,\, \text{e} \,\, \bar{f}^{(0)}_2 = -\dfrac{1}{4}.
\]
Como $f^{(n)} = V D^n V^{-1} f^{(0)} = \bar{f}^{(0)}_1 \lambda_1^n v_1 + \bar{f}^{(0)}_2 \lambda_2^n v_2$, lendo a primeira componente:
\[
  f(n) = \dfrac{5}{4} \cdot 3^n - \dfrac{1}{4} (-1)^n = \dfrac{5 \cdot 3^n - (-1)^n}{4}.
\]

O mesmo resultado sai substituindo direto na fórmula geral do capítulo da diagonalização:
\[
  f(n) = \dfrac{3^n (1 \cdot (-1) - 4) + (-1)^n (4 - 1 \cdot 3)}{-1 - 3} = \dfrac{-5 \cdot 3^n + (-1)^n}{-4} = \dfrac{5 \cdot 3^n - (-1)^n}{4}.
\]

\section{Pela função geradora}

Em notação de streams, a recorrência é
\[
  a = 1 \prec \big( 2a + (4 - 2) \prec (3a) \big).
\]
Convertendo para função geradora:
\begin{calcproof}[.]
  \start{A(x) = 1 + x \big( 2 A(x) + 2 + x \cdot 3 A(x) \big)}
  \step[\calciff]{Distributiva}{
    A(x) = 1 + 2x + 2x A(x) + 3x^2 A(x)
  }
  \step[\calciff]{Subtração dos dois lados e Distributiva}{
    A(x) (1 - 2x - 3x^2) = 1 + 2x
  }
  \step[\calciff]{Divisão dos dois lados}{
    A(x) = \dfrac{1 + 2x}{1 - 2x - 3x^2}
  }
\end{calcproof}
O numerador é $c + (d - \alpha c)x = 1 + (4 - 2)x = 1 + 2x$, como previsto pela fórmula geral.

As raízes do denominador são
\[
  R_i = \dfrac{\lambda_i}{-\beta}:
  \quad
  R_1 = \dfrac{3}{-3} = -1,
  \quad
  R_2 = \dfrac{-1}{-3} = \dfrac{1}{3},
\]
e de fato $1 - 2x - 3x^2 = (1 - 3x)(1 + x)$ se anula em $x = -1$ e em $x = 1/3$. Valem as relações do capítulo anterior:
\[
  R_1 R_2 = -\dfrac{1}{3} = -\dfrac{1}{\beta}, \,\, R_1 = \dfrac{1}{\lambda_2} \,\, \text{e} \,\, R_2 = \dfrac{1}{\lambda_1}.
\]

Fazendo as frações parciais:
\begin{calcproof}[.]
  \start{\dfrac{1 + 2x}{(1 - 3x)(1 + x)} = \dfrac{l_1}{1 - 3x} + \dfrac{l_2}{1 + x}}
  \step[\calciff]{Multiplicação pelo denominador comum}{
    1 + 2x = l_1 (1 + x) + l_2 (1 - 3x)
  }
  \step[\calciff]{Igualdade de polinômios (coeficientes de $x^0$ e $x^1$)}{
    \begin{cases}
      l_1 + l_2 = 1 \\
      l_1 - 3 l_2 = 2
    \end{cases}
  }
  \step[\calciff]{Subtração das equações}{
    l_2 = -\dfrac{1}{4} \ \land \ l_1 = \dfrac{5}{4}
  }
\end{calcproof}

Os coeficientes são as mesmas coordenadas $\bar{f}^{(0)}_1, \bar{f}^{(0)}_2$ da diagonalização. Isso não é coincidência: somando duas vezes a primeira equação à segunda, ela vira $3 l_1 - l_2 = 4$, e o sistema fica
\[
  \begin{bmatrix}
    1 & 1 \\
    3 & -1
  \end{bmatrix}
  \begin{bmatrix}
    l_1 \\
    l_2
  \end{bmatrix}
  =
  \begin{bmatrix}
    1 \\
    4
  \end{bmatrix},
\]
que é exatamente $V l = f^{(0)}$, o mesmo sistema que define $\bar{f}^{(0)}$.

Por fim, como cada fração é uma série geométrica,
\[
  A(x) = \dfrac{5}{4} \sum_{n \geq 0} 3^n x^n - \dfrac{1}{4} \sum_{n \geq 0} (-1)^n x^n
  \quad \Longrightarrow \quad
  f(n) = \dfrac{5 \cdot 3^n - (-1)^n}{4}.
\]

\section{Conferindo}

\begin{center}
  \renewcommand{\arraystretch}{1.3}
  \begin{tabular}{c|ccccccc}
    $n$ & 0 & 1 & 2 & 3 & 4 & 5 & 6 \\
    \hline
    recorrência & 1 & 4 & 11 & 34 & 101 & 304 & 911 \\
    $\dfrac{5 \cdot 3^n - (-1)^n}{4}$ & 1 & 4 & 11 & 34 & 101 & 304 & 911
  \end{tabular}
\end{center}

Os dois métodos chegam à mesma fórmula, passando pelo mesmo polinômio (espelhado), pelas mesmas raízes (invertidas) e pelo mesmo sistema linear.
